A Biologia de Sistemas nasceu de uma insatisfação produtiva. Décadas de biologia molecular
produziram um catálogo extraordinário de peças --- genes, proteínas, metabólitos, vias de
sinalização --- mas a compreensão de como essas peças se articulam para gerar vida continuava
esquiva. Listar os componentes de uma célula não é o mesmo que entender a célula, assim como
conhecer cada nota de uma partitura não é suficiente para compreender uma sinfonia.

Esta disciplina propõe uma virada de perspectiva: em vez de isolar elementos, integrá-los.
Em vez de descrever, modelar e prever. Em vez de observar, perturbar e medir. O instrumento
central dessa nova biologia é a matemática --- não como adorno formal, mas como linguagem
capaz de capturar a dinâmica, a robustez e a emergência de sistemas complexos.

Este livro reúne o conteúdo do curso \textit{Bioinformática para Biologia de Sistemas},
ministrado no Departamento de Biofísica e Farmacologia da UNESP em Botucatu. Ele não
pressupõe formação em matemática avançada, mas tampouco foge das equações quando elas são
o caminho mais direto para a clareza. O leitor encontrará aqui um percurso que vai da
estrutura molecular do DNA até a medicina personalizada, passando por redes gênicas,
cascatas de sinalização, modelos populacionais e biologia sintética.

\medskip
\textbf{Como este livro está organizado.}
A obra divide-se em quatro partes. A \textbf{Parte~I} estabelece os fundamentos: o que são
sistemas biológicos, como se constroem modelos matemáticos e o que as redes estáticas revelam
sobre a organização da vida. A \textbf{Parte~II} aprofunda as ferramentas matemáticas
indispensáveis --- equações diferenciais, análise de estabilidade, bifurcações, estimação
de parâmetros --- que permitem transformar observações qualitativas em afirmações quantitativas
rigorosas. A \textbf{Parte~III} aplica esse instrumental aos grandes sistemas biológicos:
regulação gênica, proteínas e dobramento, metabolismo, sinalização celular e dinâmica
populacional. Por fim, a \textbf{Parte~IV} explora as fronteiras do campo: análise integrada
multi-ômica, fisiologia cardíaca como sistema de controle, medicina de sistemas,
engenharia genética e os horizontes abertos pela inteligência artificial aplicada à biologia.

\medskip
Cada capítulo foi escrito para ser lido de forma contínua, como um texto --- não como uma
sequência de tópicos desconexos. As equações aparecem quando acrescentam precisão que as
palavras não conseguem dar. Os exemplos de código ilustram como os conceitos se traduzem
em análise computacional reproduzível. Os exercícios convidam o leitor a sair da posição
passiva e colocar a mão na massa.

A missão deste curso, e deste livro, pode ser resumida em uma frase: conectar dados a
mecanismos, e números a narrativas biológicas.

\bigskip
\begin{flushright}
\textit{Ney Lemke}\\
Botucatu, \the\year
\end{flushright}
